% First-public-launch fragment. Replace only the Purpose and use block
% after the title page and before \tableofcontents. Retain the surrounding
% title page, contents, notation, conventions and scientific chapters.
\chapter*{Purpose and use}
\addcontentsline{toc}{chapter}{Purpose and use}

Scientific evidence limits the explanations a system can have. Those limits
depend on how the system is prepared, which interventions are possible, what
the detector retains, and how accurately the experiment is calibrated.
\emph{Constructive Access} develops a common framework for turning these
conditions into a checkable conclusion: a property that the evidence
identifies, an experiment that separates alternatives, a resource that every
compatible explanation requires, or a precise reason to withhold judgment.

The central mathematical object is the probability law of the complete
accessible record under a permitted experiment. An \emph{access contract}
specifies the target, candidate physical models, preparation, controls,
detector, timing, calibration, error objective and resource accounting.
Models that give the same accessible laws cannot be separated within that
contract. This familiar identification principle becomes constructive when
one can compute an obstruction, synthesize an informative protocol, or bound
the target over all compatible models.

This is a unified monograph and accompanying research toolkit. It starts
with a two-state example and develops the mathematical machinery, physical
models, finite-data procedures and reproducible applications needed to use
that principle. Three questions remain distinct throughout: whether a
distinction exists in ideal evidence, whether a finite experiment can resolve
it reliably, and what physical resources the experiment requires. Exact
identification alone does not imply an efficient or realizable measurement.

\section*{The main contributions}

The general foundations draw on established system identification, automata
and realization theory, quantum instruments, statistical decision theory,
stochastic thermodynamics and concentration inequalities. Their attribution
is part of the reference. The contributions developed here are particular
proofs, constructions and their operational connection, with a claim ledger
that separates these from the inherited mathematics.

For finite classical and quantum instruments, accessible linear spaces
support exact equivalence certificates and finite distinguishing sequences.
Exported certificates permit a separate checker to verify the claimed
obstruction or additive witness charge. Calibrated testing and bounded
decision design then address the different problem of obtaining reliable
evidence. Worked intervention models distinguish current magnitude from
orientation, hidden states from hidden reaction channels, and a missing
measurement from a missing physical assumption.

The physical-resource results ask what every compatible explanation must
contain. Finite temporal observations can impose classical state requirements
under declared preparation and detector rules. Correlation evidence can
require preparation resources or sustained dissipation under different
contracts. Explicit equilibrium realizations and sampling aliases identify
where those conclusions cease to follow. A state-count obstruction is
therefore presented with its comparison class, rather than as a general
test of quantumness.

A central diffusion result goes further than detecting irreversibility in
the recorded signal. Within a specified class of ordinary smooth compact
hidden diffusions, some completely time-reversible scalar process laws have
an attained positive minimum entropy production over every compatible
realization, allowing arbitrary finite hidden dimension. For the family
\[
 L=\gamma\partial_\phi^2+c(z)\partial_\phi+L_Z,
 \qquad Y=\cos\phi,
\]
where $Z$ is a reversible compact diffusion, that minimum is
$\mathbb E[c(Z)^2]/\gamma$. It can remain positive when the mean angular
current is zero. The claim concerns the complete ideal observation law and
ordinary reversal of all hidden coordinates. Changing the admissible dynamics
or the physical reversal convention changes the question.

The detector and calibration results connect such population statements to
finite observations. A replicated polynomial construction supplies an entropy
lower confidence bound when the bounded observable has the pointwise geometry
$\Gamma_f=\gamma(1-f^2)$ with an unknown scale $\gamma$. Temporal moments,
constant references and timed-step references cover the scale, a common
detector gain and a uniform exposure width. The geometric relation and the
required generator bounds are explicit premises. A complementary rare-region
theorem proves why finite passive records cannot uniformly establish
unrestricted global bounds merely from individual-model smoothness, favorable
ellipticity or a mixing gap.

The sharp polynomial limit is constructive but requires increasing access,
shrinking exposure and increasingly accurate arithmetic. Finite performance
is reported separately. In the frozen calibration experiment, each dataset
uses $13{,}333{,}333$ independent target trials and $200{,}000$ reference
readings, totaling $80{,}199{,}998$ detector readings. Positive geometry
certificates occur in the stronger-drive cases; the weak-drive case remains
inconclusive, and the higher polynomial prefix never wins at that budget.
The broader-class endpoint comparator performs better in some detection
comparisons. These outcomes expose remaining acquisition costs rather than
establishing universal efficiency. Detector counts and elapsed time are
reported separately from energy whenever apparatus energy has not been
calibrated.

\section*{Reading and using the reference}

The chapters are arranged by dependency, with several entry routes:
\begin{itemize}[leftmargin=1.5em]
\item \textbf{First reading.}
Chapter~\ref{entry:start} introduces the access contract through a complete
example. Chapters~\ref{foundations:chapter}, \ref{finite:chapter} and
\ref{precision:chapter} develop identification, constructive certificates
and finite-error inference. Chapter~\ref{workflow:suite} assembles these
into an experimental workflow.
\item \textbf{Controlled classical and quantum experiments.}
Begin with Chapter~\ref{finite:chapter}, then use the response, preparation,
dimension and application chapters for the appropriate physical contract.
Check which controls and retained outcomes a proposed instrument permits
before applying its certificate.
\item \textbf{Dissipation and hidden diffusion.}
Chapter~\ref{pa:entry} introduces correlation evidence for necessary
dissipation. Chapters~\ref{og:entry}, \ref{ch:dg-entry} and
\ref{ch:dc-entry} connect observed diffusion geometry to physical detectors
and uncertain calibration. Read each constructive result together with its
countermodels and acquisition assumptions.
\item \textbf{Implementation and reproduction.}
Chapter~\ref{toolkit:chapter} gives installation instructions and examples
for the \texttt{constructive\_access} package. The archive provides
proof sources, frozen designs, retained statistics, provenance and verification
drivers. The calibration tutorial replays a retained simulated dataset through
the public exact-arithmetic API. Chapter~\ref{vl:entry} explains what the
virtual experiments establish.
\end{itemize}

A practical application begins by matching an apparatus to a contract, then
choosing the target and the corresponding certificate. Inconclusive output is
an informative outcome when the remaining models disagree or the confidence
region is too wide. Additional sensing, intervention, preparation or calibration
may resolve it; additional processing of the same indistinguishable evidence
cannot. The software checks the stated calculation. Physical calibration and
model adequacy must be justified in the application.

\section*{Evidence and review}

The evidential basis comprises mathematical arguments, internally checked
software, synthetic experiments and an explicitly exploratory reanalysis of
previously published optical-tweezer measurements. The latter identifies a
calibration gap and supplies no empirical entropy certificate. This work
reports no newly measured apparatus or experimental verification of a new
physical law. The claim ledger records assumptions and evidence status, and
the literature chapters distinguish established methods from the specific
results presented here. Historical priority has not been established.

The monograph has not undergone external peer review. The
released proofs, sources and reproduction records are provided so that readers
can inspect, challenge and extend the results under their stated assumptions.

\section*{Availability and reuse}

Project information and correspondence:
\url{https://2dlabs.org} and
\href{mailto:helloworld@2dai.org}{\nolinkurl{helloworld@2dai.org}}.
Copyright \textcopyright{} 2026 Abdinoor (2D Labs), to the extent of rights
owned or controlled by the licensor. The covered research text and software
use the \emph{Constructive Access Research and Revenue-Threshold License},
version~1.0. Its controlling terms accompany the research package as
\texttt{LICENSE}, with a \texttt{COMMERCIAL\_USE.txt} guide; copies and
commercial terms may be requested at the address above.

Personal noncommercial use and qualifying academic research and teaching are
free. Business use is free while worldwide Group Revenue, including controlled
and common-control affiliates, is at or below US\$10 million over the preceding
12 months, subject to the license's eligibility conditions. Above that
threshold, covered business use outside the universally permitted scholarly
activities requires a separate written commercial license. Reading, scientific
review, reproduction of the published demonstrations and scholarly quotation
remain permitted under the stated terms regardless of revenue. The license
does not reserve mathematical truths, scientific ideas or independent
implementations that copy no protected expression. Third-party material
retains its own licenses.
